\label{curvilinear}In the following, we note ${\bf X}=\left( x,y,z\right) $
the vector position. In general, vectors are written in bold characters,
except unit vectors, noted with a hat. We consider the following curvilinear
coordinate systems

\subsection{General Case $(u^{1},u^{2},u^{3})$}

We consider the curvilinear coordinate system $(u^{1},u^{2},u^{3})$.

\subsubsection{Covariant (Tangent) Basis}

The covariant, or tangent vector basis $({\bf e}_{1},{\bf e}_{2},{\bf e}%
_{3}) $is defined as 
\begin{equation}
{\bf e}_{i}=\frac{\partial {\bf X}}{\partial u^{i}}  \label{cov}
\end{equation}
where the ${\bf e}_{i}$ are tangent to the curvilinear lines. They can be
normalized by 
\begin{equation}
\widehat{e}_{i}=\frac{{\bf e}_{i}}{h_{i}}
\end{equation}
where we introduce the scale factors 
\begin{equation}
h_{i}=\left| \frac{\partial {\bf X}}{\partial u^{i}}\right|
\end{equation}

\subsubsection{Contravariant (reciprocal) Basis}

The gradient $\nabla f$ of a function $f$ being defined by the differential 
\begin{equation}
df=\nabla f\cdot d{\bf X}
\end{equation}
we apply to $u^{i}$ which gives 
\begin{equation}
du^{i}=\nabla u^{i}\cdot d{\bf X}
\end{equation}

By chain rule, we have 
\begin{equation}
d{\bf X=}\frac{\partial {\bf X}}{\partial u^{j}}du^{j}={\bf e}_{j}du^{j}
\label{dX}
\end{equation}
so that 
\begin{equation}
du^{i}=\nabla u^{i}\cdot {\bf e}_{j}du^{j}
\end{equation}
which implies 
\begin{equation}
\nabla u^{i}\cdot {\bf e}_{j}=\delta _{j}^{i}
\end{equation}
thus defining two reciprocal basis $\left( \nabla u^{i},{\bf e}_{j}\right) $%
of vectors. The reciprocal basis vectors are also called contravariant, and
noted 
\begin{equation}
{\bf e}^{i}\equiv \nabla u^{i}  \label{cont}
\end{equation}
These vectors are perpendicular to the surfaces of constant $u^{i}$.

From the properties of reciprocal basis, we can calculate a vector from the
three vectors of the reciprocal basis, such that 
\begin{eqnarray}
{\bf e}^{i} &=&\nabla u^{i}=\frac{{\bf e}_{j}\times {\bf e}_{k}}{{\bf e}%
_{i}\cdot {\bf e}_{j}\times {\bf e}_{k}}  \label{rel1} \\
{\bf e}_{i} &=&\frac{\partial {\bf X}}{\partial u^{i}}=\frac{{\bf e}%
^{j}\times {\bf e}^{k}}{{\bf e}^{i}\cdot {\bf e}^{j}\times {\bf e}^{k}}
\label{rel2}
\end{eqnarray}

\subsubsection{Metric Coefficients}

They are defined as 
\begin{eqnarray}
g_{ij} &=&{\bf e}_{i}\cdot {\bf e}_{j}  \label{met} \\
g^{ij} &=&{\bf e}^{i}\cdot {\bf e}^{j}  \nonumber
\end{eqnarray}

With the differential vector given in (\ref{dX}), we show that the
differential arc length along a curve is 
\begin{equation}
dl=\left| d{\bf X}\right| {\bf =}\sqrt{d{\bf X\cdot }d{\bf X}}=\sqrt{%
g_{ij}du^{i}du^{j}}  \label{dl}
\end{equation}

In addition, we have the relations 
\begin{eqnarray}
{\bf e}_{i} &=&g_{ij}{\bf e}^{j}  \label{corr} \\
{\bf e}^{i} &=&g^{ij}{\bf e}_{j}  \nonumber
\end{eqnarray}

We also show that 
\begin{eqnarray}
\left[ g_{ij}\right] &=&\left[ g^{ij}\right] ^{-1} \\
\left[ g^{ij}\right] &=&\left[ g_{ij}\right] ^{-1}
\end{eqnarray}
so that, defining 
\begin{equation}
g=\det \left[ g_{ij}\right]
\end{equation}
we find 
\begin{equation}
g^{-1}=\det \left[ g^{ij}\right]
\end{equation}

\subsubsection{Jacobian}

We define the Jacobian 
\begin{equation}
J\equiv \frac{\partial (x,y,z)}{\partial (u^{1},u^{2},u^{3})}=\det \left( 
\begin{array}{lll}
\partial x/\partial u^{1} & \partial x/\partial u^{2} & \partial x/\partial
u^{3} \\ 
\partial y/\partial u^{1} & \partial y/\partial u^{2} & \partial y/\partial
u^{3} \\ 
\partial z/\partial u^{1} & \partial z/\partial u^{2} & \partial z/\partial
u^{3}
\end{array}
\right)
\end{equation}
which gives 
\begin{equation}
J=\frac{\partial {\bf X}}{\partial u^{1}}\cdot \frac{\partial {\bf X}}{%
\partial u^{2}}\times \frac{\partial {\bf X}}{\partial u^{3}}={\bf e}%
_{1}\cdot {\bf e}_{2}\times {\bf e}_{3}
\end{equation}
and the reciprocal Jacobian 
\begin{equation}
{\cal J}\equiv \frac{\partial (u^{1},u^{2},u^{3})}{\partial (x,y,z)}=\det
\left( 
\begin{array}{lll}
\partial u^{1}/\partial x & \partial u^{1}/\partial y & \partial
u^{1}/\partial z \\ 
\partial u^{2}/\partial x & \partial u^{2}/\partial y & \partial
u^{2}/\partial z \\ 
\partial u^{3}/\partial x & \partial u^{3}/\partial y & \partial
u^{3}/\partial z
\end{array}
\right)
\end{equation}
which gives 
\begin{equation}
{\cal J}=\nabla u^{1}\cdot \nabla u^{2}\times \nabla u^{3}={\bf e}^{1}\cdot 
{\bf e}^{2}\times {\bf e}^{3}
\end{equation}

We can show that 
\begin{equation}
{\cal J}=J^{-1}
\end{equation}
and the relations (\ref{rel1}-\ref{rel2}) become 
\begin{eqnarray}
{\bf e}^{i} &=&\frac{1}{J}\left( {\bf e}_{j}\times {\bf e}_{k}\right) \\
{\bf e}_{i} &=&J\left( {\bf e}^{j}\times {\bf e}^{k}\right)
\end{eqnarray}

Also, 
\begin{equation}
g=J^{2}
\end{equation}

\subsubsection{Vector Identities}

With 
\begin{eqnarray}
{\bf A} &=&\left( {\bf A}\cdot {\bf e}_{i}\right) {\bf e}^{i}=A_{i}{\bf e}%
^{i} \\
{\bf A} &=&\left( {\bf A}\cdot {\bf e}^{i}\right) {\bf e}_{i}=A^{i}{\bf e}%
_{i}
\end{eqnarray}
we find 
\begin{equation}
{\bf A}\cdot {\bf B}=g_{ij}A^{i}B^{j}=g^{ij}A_{i}B_{j}
\end{equation}
so that 
\begin{equation}
A=\left| {\bf A}\right| =\sqrt{g_{ij}A^{i}A^{j}}=\sqrt{g^{ij}A_{i}A_{j}}
\end{equation}

We also find 
\begin{equation}
{\bf A}\times {\bf B}=A^{i}B^{j}{\bf e}_{i}\times {\bf e}_{j}=A_{i}B_{j}{\bf %
e}^{i}\times {\bf e}^{j}
\end{equation}
which gives 
\begin{equation}
\left( {\bf A}\times {\bf B}\right) _{k}=\varepsilon _{ijk}JA^{i}B^{j}=\frac{%
\varepsilon ^{ijk}}{J}A_{i}B_{j}
\end{equation}

Note that from (\ref{corr}), 
\begin{eqnarray}
A_{i} &=&g_{ij}A^{j} \\
A^{i} &=&g^{ij}A_{j}
\end{eqnarray}

\subsubsection{Differential elements}

\paragraph{differential length along $u^{i}$}

\begin{equation}
dl\left( i\right) =\left| d{\bf X}\left( i\right) \right| =h_{i}du^{i}=\sqrt{%
g_{ii}}du^{i}
\end{equation}

Equivalently, 
\begin{equation}
dl\left( i\right) =J\left| \nabla u^{j}\times \nabla u^{k}\right| du^{i}
\end{equation}

\paragraph{differential area in surface of constant $u^{i}$}

Using 
\begin{equation}
dS\left( i\right) =\left| d{\bf X}\left( j\right) \times d{\bf X}\left(
k\right) \right| =\left| {\bf e}_{j}\times {\bf e}_{k}\right| du^{j}du^{k}
\end{equation}
which becomes 
\begin{equation}
dS\left( i\right) =\sqrt{g_{jj}g_{kk}-g_{jk}^{2}}du^{j}du^{k}
\end{equation}

Equivalently 
\begin{equation}
dS\left( i\right) =J\left| \nabla u^{i}\right| du^{j}du^{k}
\end{equation}
so that 
\begin{equation}
d{\bf S}\left( i\right) =\pm Jdu^{j}du^{k}\nabla u^{i}
\end{equation}

\paragraph{differential volume element}

\begin{equation}
d^{3}{\bf X=}d{\bf X}\left( 1\right) \cdot d{\bf X}\left( 2\right) \times d%
{\bf X}\left( 3\right) =Jdu^{1}du^{2}du^{3}
\end{equation}

\subsubsection{Vector Differentiation}

\begin{equation}
d{\bf A}=\frac{\partial {\bf A}}{\partial u^{k}}du^{k}
\end{equation}

\begin{eqnarray}
\frac{\partial {\bf A}}{\partial u^{k}} &=&\left( \frac{\partial {\bf A}}{%
\partial u^{k}}\right) ^{j}{\bf e}_{j}\equiv {\bf A}_{k}^{j}{\bf e}_{j} \\
&=&\left( \frac{\partial {\bf A}}{\partial u^{k}}\right) _{j}{\bf e}%
^{j}\equiv {\bf A}_{j,k}{\bf e}^{j}
\end{eqnarray}
with 
\begin{eqnarray}
{\bf A}_{k}^{j} &=&\frac{\partial A^{j}}{\partial u^{k}}+A^{i}\left\{ \frac{%
\partial {\bf e}_{i}}{\partial u^{k}}\cdot {\bf e}^{j}\right\} \\
{\bf A}_{j,k} &=&\frac{\partial A_{j}}{\partial u^{k}}+A_{i}\left\{ \frac{%
\partial {\bf e}^{i}}{\partial u^{k}}\cdot {\bf e}_{j}\right\}
\end{eqnarray}

Then, 
\begin{eqnarray}
\delta {\bf A}^{j} &=&\left( d{\bf A}\right) ^{j}=\frac{\partial A^{j}}{%
\partial u^{k}}du^{k}+A^{i}\left\{ \frac{\partial {\bf e}_{i}}{\partial u^{k}%
}\cdot {\bf e}^{j}\right\} du^{k} \\
&=&dA^{j}+\left\{ \frac{\partial {\bf e}_{i}}{\partial u^{k}}\cdot {\bf e}%
^{j}\right\} A^{i}du^{k}
\end{eqnarray}

It can be shown that the Christoffel Symbol of the second kind is 
\begin{equation}
\left\{ \frac{\partial {\bf e}_{i}}{\partial u^{k}}\cdot {\bf e}^{j}\right\}
\equiv \left\{ 
\begin{array}{l}
j \\ 
\begin{array}{ll}
i & k
\end{array}
\end{array}
\right\} =\frac{1}{2}g^{jn}\left[ \frac{\partial g_{ni}}{\partial u^{k}}+%
\frac{\partial g_{nk}}{\partial u^{i}}-\frac{\partial g_{ik}}{\partial u^{n}}%
\right]  \label{chris}
\end{equation}
so that 
\begin{eqnarray}
{\bf A}_{k}^{j} &=&\frac{\partial A^{j}}{\partial u^{k}}+\left\{ 
\begin{array}{l}
j \\ 
\begin{array}{ll}
i & k
\end{array}
\end{array}
\right\} A^{i} \\
{\bf A}_{j,k} &=&\frac{\partial A_{j}}{\partial u^{k}}-\left\{ 
\begin{array}{l}
i \\ 
\begin{array}{ll}
j & k
\end{array}
\end{array}
\right\} A_{i}
\end{eqnarray}

Note that since 
\begin{equation}
\frac{\partial {\bf e}_{i}}{\partial u^{k}}=\frac{\partial {\bf e}_{k}}{%
\partial u^{i}}
\end{equation}
we have 
\begin{equation}
\left\{ 
\begin{array}{l}
j \\ 
\begin{array}{ll}
i & k
\end{array}
\end{array}
\right\} =\left\{ 
\begin{array}{l}
j \\ 
\begin{array}{ll}
k & i
\end{array}
\end{array}
\right\}
\end{equation}

\subsubsection{Operator $\nabla $}

The operator $\nabla $ can be decomposed in the curvilinear coordinates as 
\begin{equation}
\nabla =\nabla u^{i}\frac{\partial }{\partial u^{i}}={\bf e}^{i}\frac{%
\partial }{\partial u^{i}}
\end{equation}

We then find the following differential operations:

\paragraph{Gradient}

It follows simply that 
\begin{equation}
\nabla f=\nabla u^{i}\frac{\partial f}{\partial u^{i}}=\frac{\partial f}{%
\partial u^{i}}{\bf e}^{i}
\end{equation}
so that 
\begin{equation}
\left( \nabla f\right) _{i}=\left( \nabla f\cdot {\bf e}_{i}\right) =\frac{%
\partial f}{\partial u^{i}}
\end{equation}

\paragraph{Divergence}

It can be shown that the divergence is expressed as 
\begin{equation}
\nabla \cdot {\bf A}=\frac{1}{J}\frac{\partial }{\partial u^{i}}\left(
JA^{i}\right)
\end{equation}

\paragraph{Curl}

It becomes, a compact notations, 
\begin{equation}
\nabla \times {\bf A}=\frac{\varepsilon ^{ijk}}{J}\frac{\partial A_{j}}{%
\partial u^{i}}{\bf e}_{k}
\end{equation}
or is extended as 
\begin{equation}
\left( \nabla \times {\bf A}\right) ^{k}=\frac{1}{J}\left( \frac{\partial
A_{j}}{\partial u^{i}}-\frac{\partial A_{i}}{\partial u^{j}}\right)
\end{equation}

\subsection{System $\left( R,Z,\phi \right) $}

\subsubsection{Definition}

The coordinates $\left( R,Z,\phi \right) $ are defined on the space 
\begin{eqnarray}
0\leq R &<&\infty  \nonumber \\
-\infty \leq Z &<&\infty  \nonumber \\
0\leq \phi &<&2\pi
\end{eqnarray}
and is related to $\left( x,y,z\right) $ by

\begin{eqnarray}
R &=&\sqrt{x^{2}+y^{2}}  \label{RZphidef} \\
Z &=&-z  \nonumber \\
\phi &=&\arctan \left( y/x\right) +\pi H\left( -x\right) \;\ \left[ 2\pi
\right]  \nonumber
\end{eqnarray}
which is inverted to 
\begin{eqnarray}
x &=&R\cos \phi \\
y &=&R\sin \phi  \nonumber \\
z &=&-Z  \nonumber
\end{eqnarray}

\subsubsection{Position Vector}

The position vector then becomes 
\begin{equation}
{\bf X}=R\widehat{R}+Z\widehat{Z}
\end{equation}
where we define a local orthonormal basis $\left( \widehat{R},\widehat{Z},%
\widehat{\phi }\right) $ as 
\begin{eqnarray}
\widehat{R} &=&\cos \phi \;\widehat{x}+\sin \phi \;\widehat{y} \\
\widehat{Z} &=&-\widehat{z} \\
\widehat{\phi } &=&\widehat{R}\times \widehat{Z}=-\sin \phi \;\widehat{x}%
+\cos \phi \;\widehat{y}
\end{eqnarray}

\subsubsection{Covariant Basis}

The covariant vector basis is defined in (\ref{cov}), which becomes here 
\begin{eqnarray}
{\bf e}_{R} &=&\frac{\partial {\bf X}}{\partial R}=\widehat{R} \\
{\bf e}_{Z} &=&\frac{\partial {\bf X}}{\partial Z}=\widehat{Z} \\
{\bf e}_{\phi } &=&\frac{\partial {\bf X}}{\partial \phi }=R\frac{\partial 
\widehat{R}}{\partial \phi }=R\widehat{\phi }
\end{eqnarray}
so that we have the covariant basis 
\begin{equation}
\left( {\bf e}_{R},{\bf e}_{Z},{\bf e}_{\phi }\right) =\left( \widehat{R},%
\widehat{Z},R\widehat{\phi }\right)
\end{equation}
the scaling factors 
\begin{equation}
\left( h_{R},h_{Z},h_{\phi }\right) =\left( 1,1,R\right)
\end{equation}
and the normalized tangent basis 
\begin{equation}
\left( \widehat{e}_{R},\widehat{e}_{Z},\widehat{e}_{\phi }\right) =\left( 
\widehat{R},\widehat{Z},\widehat{\phi }\right)
\end{equation}

\subsubsection{Contravariant Basis}

The Contravariant vector basis is defined in (\ref{cont}), which becomes
here 
\begin{eqnarray}
{\bf e}^{R} &=&\nabla R=\widehat{R} \\
{\bf e}^{Z} &=&\nabla Z=\widehat{Z} \\
{\bf e}^{\phi } &=&\nabla \phi =\frac{\widehat{\phi }}{R}
\end{eqnarray}

The relations (\ref{rel1}-\ref{rel2}) are here readily verified. The
normalized reciprocal basis is 
\begin{equation}
\left( \widehat{e}^{R},\widehat{e}^{Z},\widehat{e}^{\phi }\right) =\left( 
\widehat{R},\widehat{Z},\widehat{\phi }\right)
\end{equation}
which here coincides with the normalized tangent basis, since both bases are
orthogonal.

\subsubsection{Metric Coefficients}

They are defined in (\ref{met}) and become here 
\begin{eqnarray}
g_{ij} &=&\left( 
\begin{array}{lll}
1 & 0 & 0 \\ 
0 & 1 & 0 \\ 
0 & 0 & R^{2}
\end{array}
\right) \\
g^{ij} &=&\left( 
\begin{array}{lll}
1 & 0 & 0 \\ 
0 & 1 & 0 \\ 
0 & 0 & 1/R^{2}
\end{array}
\right)  \nonumber
\end{eqnarray}

As a result 
\begin{equation}
g=R^{2}
\end{equation}
and the Jacobian is 
\begin{equation}
J=R
\end{equation}

\subsubsection{Differential elements}

\begin{eqnarray}
dl\left( R\right) &=&dR \\
dl\left( Z\right) &=&dZ \\
dl\left( \phi \right) &=&Rd\phi
\end{eqnarray}

\begin{eqnarray}
d{\bf S}\left( R\right) &=&RdZd\phi \widehat{R} \\
d{\bf S}\left( Z\right) &=&RdRd\phi \widehat{Z} \\
d{\bf S}\left( \phi \right) &=&dRdZ\widehat{\phi }
\end{eqnarray}

\begin{equation}
d^{3}{\bf X}=RdRdZd\phi
\end{equation}

\subsubsection{Christoffel Symbols}

They are defined in (\ref{chris}) and are all zero here except

\[
\left\{ 
\begin{array}{l}
Z \\ 
\begin{array}{ll}
Z & R
\end{array}
\end{array}
\right\} =\left\{ 
\begin{array}{l}
Z \\ 
\begin{array}{ll}
R & Z
\end{array}
\end{array}
\right\} =\frac{1}{2}g^{ZZ}\frac{\partial g_{ZZ}}{r}=\frac{1}{R} 
\]
\begin{equation}
\left\{ 
\begin{array}{l}
R \\ 
\begin{array}{ll}
Z & Z
\end{array}
\end{array}
\right\} =-\frac{1}{2}g^{RR}\frac{\partial g_{ZZ}}{r}=-R
\end{equation}

\subsubsection{Differential Operations}

\paragraph{Gradient}

\begin{equation}
\nabla f=\frac{\partial f}{\partial R}\widehat{R}+\frac{\partial f}{\partial
Z}\widehat{Z}+\frac{1}{R}\frac{\partial f}{\partial \phi }\widehat{\phi }
\end{equation}

\paragraph{Divergence}

\begin{equation}
\nabla \cdot {\bf A}=\frac{1}{R}\frac{\partial }{\partial R}\left( R{\bf A}%
\cdot \widehat{R}\right) +\frac{\partial }{\partial Z}\left( {\bf A}\cdot 
\widehat{Z}\right) +\frac{1}{R}\frac{\partial }{\partial \phi }\left( {\bf A}%
\cdot \widehat{\phi }\right)
\end{equation}

\paragraph{Curl}

\begin{eqnarray}
\left( \nabla \times {\bf A}\right) \cdot \widehat{R} &=&\frac{\partial }{%
\partial Z}\left( {\bf A}\cdot \widehat{\phi }\right) -\frac{1}{R}\frac{%
\partial }{\partial \phi }\left( {\bf A}\cdot \widehat{Z}\right) \\
\left( \nabla \times {\bf A}\right) \cdot \widehat{Z} &=&\frac{1}{R}\frac{%
\partial }{\partial \phi }\left( {\bf A}\cdot \widehat{R}\right) -\frac{1}{R}%
\frac{\partial }{\partial R}\left( R{\bf A}\cdot \widehat{\phi }\right) \\
\left( \nabla \times {\bf A}\right) \cdot \widehat{\phi } &=&\frac{\partial 
}{\partial R}\left( {\bf A}\cdot \widehat{Z}\right) -\frac{\partial }{%
\partial Z}\left( {\bf A}\cdot \widehat{R}\right)
\end{eqnarray}

\subsection{System $\left( r,\theta ,\phi \right) $}

\subsubsection{Definition}

The coordinates $\left( r,\theta ,\phi \right) $ are defined from the origin 
$\left( R_{0},Z_{0}\right) $ on the space 
\begin{eqnarray}
0\leq r &<&\infty \\
0\leq \theta &<&2\pi  \nonumber
\end{eqnarray}
and is related to $\left( R,Z,\phi \right) $ by

\begin{eqnarray}
r &=&\sqrt{\left( R-R_{0}\right) ^{2}+\left( Z-Z_{0}\right) ^{2}}
\label{rthetadef} \\
\theta &=&\arctan \left( \left( Z-Z_{0}\right) /\left( R-R_{0}\right)
\right) +\pi H\left( R_{0}-R\right) \;\ \left[ 2\pi \right]  \nonumber
\end{eqnarray}
which is inverted to 
\begin{eqnarray}
R &=&R_{0}+r\cos \theta \\
Z &=&Z_{0}+r\cos \theta  \nonumber
\end{eqnarray}

\subsubsection{Position Vector}

The position vector then becomes 
\begin{equation}
{\bf X}=R_{0}\widehat{R}+Z_{0}\widehat{Z}+r\widehat{r}
\end{equation}
where we define a local orthonormal basis $\left( \widehat{r},\widehat{%
\theta },\widehat{\phi }\right) $ as 
\begin{eqnarray}
\widehat{r} &=&\cos \theta \;\widehat{R}+\sin \theta \;\widehat{Z} \\
\widehat{\theta } &=&\widehat{\phi }\times \widehat{r}=-\sin \theta \;%
\widehat{R}+\cos \theta \;\widehat{Z}
\end{eqnarray}
since 
\begin{eqnarray}
\widehat{\phi }\times \widehat{r} &=&\left( \widehat{R}\times \widehat{Z}%
\right) \times \left( \cos \theta \;\widehat{R}+\sin \theta \;\widehat{Z}%
\right) \\
&=&\left[ \left( \cos \theta \;\widehat{R}+\sin \theta \;\widehat{Z}\right)
\cdot \widehat{R}\right] \widehat{Z}-\left[ \left( \cos \theta \;\widehat{R}%
+\sin \theta \;\widehat{Z}\right) \cdot \widehat{Z}\right] \widehat{R} \\
&=&\cos \theta \widehat{Z}-\sin \theta \;\widehat{R}
\end{eqnarray}

\subsubsection{Covariant Basis}

The covariant vector basis is defined in (\ref{cov}), which becomes here 
\begin{eqnarray}
{\bf e}_{r} &=&\frac{\partial {\bf X}}{\partial r}=\widehat{r} \\
{\bf e}_{\theta } &=&\frac{\partial {\bf X}}{\partial \theta }=r\frac{%
\partial \widehat{r}}{\partial \theta }=r\widehat{\theta } \\
{\bf e}_{\phi } &=&\frac{\partial {\bf X}}{\partial \phi }=R_{0}\frac{%
\partial \widehat{R}}{\partial \phi }+r\frac{\partial \widehat{r}}{\partial
\phi }=\left( R_{0}+r\cos \theta \right) \frac{\partial \widehat{R}}{%
\partial \phi }=R\widehat{\phi }
\end{eqnarray}
so that we have the covariant basis 
\begin{equation}
\left( {\bf e}_{r},{\bf e}_{\theta },{\bf e}_{\phi }\right) =\left( \widehat{%
r},r\widehat{\theta },R\widehat{\phi }\right)
\end{equation}
the scaling factors 
\begin{equation}
\left( h_{r},h_{\theta },h_{\phi }\right) =\left( 1,r,R\right)
\end{equation}
and the normalized tangent basis 
\begin{equation}
\left( \widehat{e}_{r},\widehat{e}_{\theta },\widehat{e}_{\phi }\right)
=\left( \widehat{r},\widehat{\theta },\widehat{\phi }\right)
\end{equation}

\subsubsection{Contravariant Basis}

The Contravariant vector basis is defined in (\ref{cont}), which becomes
here 
\begin{eqnarray}
{\bf e}^{r} &=&\nabla r=\widehat{r} \\
{\bf e}^{\theta } &=&\nabla \theta =\frac{1}{r}\widehat{\theta } \\
{\bf e}^{\phi } &=&\nabla \phi =\frac{\widehat{\phi }}{R}
\end{eqnarray}

The relations (\ref{rel1}-\ref{rel2}) are here readily verified. The
normalized reciprocal basis is 
\begin{equation}
\left( \widehat{e}^{r},\widehat{e}^{\theta },\widehat{e}^{\phi }\right)
=\left( \widehat{r},\widehat{\theta },\widehat{\phi }\right)
\end{equation}
which here coincides with the normalized tangent basis, since both bases are
orthogonal.

\subsubsection{Metric Coefficients}

They are defined in (\ref{met}) and become here 
\begin{eqnarray}
g_{ij} &=&\left( 
\begin{array}{lll}
1 & 0 & 0 \\ 
0 & r^{2} & 0 \\ 
0 & 0 & R^{2}
\end{array}
\right)   \label{gijrtheta} \\
g^{ij} &=&\left( 
\begin{array}{lll}
1 & 0 & 0 \\ 
0 & 1/r^{2} & 0 \\ 
0 & 0 & 1/R^{2}
\end{array}
\right)   \nonumber
\end{eqnarray}

As a result 
\begin{equation}
g=r^{2}R^{2}
\end{equation}
and the Jacobian is 
\begin{equation}
J=rR
\end{equation}

\subsubsection{Differential elements}

\begin{eqnarray}
dl\left( r\right) &=&dr \\
dl\left( \theta \right) &=&d\theta \\
dl\left( \phi \right) &=&Rd\phi
\end{eqnarray}

\begin{eqnarray}
d{\bf S}\left( r\right) &=&rRd\theta d\phi \widehat{r} \\
d{\bf S}\left( \theta \right) &=&Rdrd\phi \widehat{\theta } \\
d{\bf S}\left( \phi \right) &=&rdrd\theta \widehat{\phi }
\end{eqnarray}

\begin{equation}
d^{3}{\bf X}=rRdrd\theta d\phi
\end{equation}

\subsubsection{Christoffel Symbols}

They are defined in (\ref{chris}) and are all zero here except 
\begin{equation}
\left\{ 
\begin{array}{l}
\theta \\ 
\begin{array}{ll}
\theta & r
\end{array}
\end{array}
\right\} =\left\{ 
\begin{array}{l}
\theta \\ 
\begin{array}{ll}
r & \theta
\end{array}
\end{array}
\right\} =\frac{1}{2}g^{\theta \theta }\frac{\partial g_{\theta \theta }}{r}=%
\frac{1}{r}
\end{equation}
\begin{equation}
\left\{ 
\begin{array}{l}
r \\ 
\begin{array}{ll}
\theta & \theta
\end{array}
\end{array}
\right\} =-\frac{1}{2}g^{rr}\frac{\partial g_{\theta \theta }}{r}=-r
\end{equation}
\begin{equation}
\left\{ 
\begin{array}{l}
\phi \\ 
\begin{array}{ll}
\phi & r
\end{array}
\end{array}
\right\} =\left\{ 
\begin{array}{l}
\phi \\ 
\begin{array}{ll}
r & \phi
\end{array}
\end{array}
\right\} =\frac{1}{2}g^{\phi \phi }\frac{\partial g_{\phi \phi }}{r}=\frac{1%
}{R}\cos \theta
\end{equation}
\begin{equation}
\left\{ 
\begin{array}{l}
r \\ 
\begin{array}{ll}
\phi & \phi
\end{array}
\end{array}
\right\} =-\frac{1}{2}g^{rr}\frac{\partial g_{\phi \phi }}{r}=-R\cos \theta
\end{equation}
\begin{equation}
\left\{ 
\begin{array}{l}
\phi \\ 
\begin{array}{ll}
\phi & \theta
\end{array}
\end{array}
\right\} =\left\{ 
\begin{array}{l}
\phi \\ 
\begin{array}{ll}
\theta & \phi
\end{array}
\end{array}
\right\} =\frac{1}{2}g^{\phi \phi }\frac{\partial g_{\phi \phi }}{\theta }=-%
\frac{r}{R}\sin \theta
\end{equation}
\begin{equation}
\left\{ 
\begin{array}{l}
\theta \\ 
\begin{array}{ll}
\phi & \phi
\end{array}
\end{array}
\right\} =-\frac{1}{2}g^{\theta \theta }\frac{\partial g_{\phi \phi }}{%
\theta }=\frac{R}{r}\sin \theta
\end{equation}

\subsubsection{Differential Operations}

\paragraph{Gradient}

\begin{equation}
\nabla f=\frac{\partial f}{\partial r}\widehat{r}+\frac{1}{r}\frac{\partial f%
}{\partial \theta }\widehat{\theta }+\frac{1}{R}\frac{\partial f}{\partial
\phi }\widehat{\phi }
\end{equation}

\paragraph{Divergence}

\begin{equation}
\nabla \cdot {\bf A}=\frac{1}{rR}\frac{\partial }{\partial r}\left( rR{\bf A}%
\cdot \widehat{r}\right) +\frac{1}{rR}\frac{\partial }{\partial \theta }%
\left( R{\bf A}\cdot \widehat{\theta }\right) +\frac{1}{R}\frac{\partial }{%
\partial \phi }\left( {\bf A}\cdot \widehat{\phi }\right)
\end{equation}

\paragraph{Curl}

\begin{eqnarray}
\left( \nabla \times {\bf A}\right) \cdot \widehat{r} &=&\frac{1}{rR}\frac{%
\partial }{\partial \theta }\left( R{\bf A}\cdot \widehat{\phi }\right) -%
\frac{1}{R}\frac{\partial }{\partial \phi }\left( {\bf A}\cdot \widehat{%
\theta }\right) \\
\left( \nabla \times {\bf A}\right) \cdot \widehat{\theta } &=&\frac{1}{R}%
\frac{\partial }{\partial \phi }\left( {\bf A}\cdot \widehat{r}\right) -%
\frac{1}{R}\frac{\partial }{\partial r}\left( R{\bf A}\cdot \widehat{\phi }%
\right) \\
\left( \nabla \times {\bf A}\right) \cdot \widehat{\phi } &=&\frac{1}{r}%
\frac{\partial }{\partial r}\left( r{\bf A}\cdot \widehat{\theta }\right) -%
\frac{1}{r}\frac{\partial }{\partial \theta }\left( {\bf A}\cdot \widehat{r}%
\right)
\end{eqnarray}

\subsection{System $\left( \psi ,s,\phi \right) $}

\subsubsection{Definition}

The coordinates $\left( \psi ,s,\phi \right) $, used to parametrize closed
flux-surfaces, are defined from the origin $\left( R_{0},Z_{0}\right) $ on
the (closed) space 
\begin{eqnarray}
\min \left( \psi _{0},\psi _{a}\right) &\leq &\psi \leq \max \left( \psi
_{0},\psi _{a}\right) \\
&&0\leq s\leq s_{\max }
\end{eqnarray}
and is related to $\left( r,\theta ,\phi \right) $ by

\begin{eqnarray}
\psi &=&\psi \left( r,\theta \right)  \label{psisdef} \\
s &=&s\left( r,\theta \right)  \nonumber
\end{eqnarray}
which is inverted to 
\begin{eqnarray*}
r &=&r\left( \psi ,s\right) \\
\theta &=&\theta \left( \psi ,s\right)
\end{eqnarray*}

Note that $\psi \left( r,\theta \right) $ must be a monotonic function of $r$
from $\psi _{0}$ at the center $\left( R_{0},Z_{0}\right) $ to $\psi _{a}$
at the edge. It is the case for nested flux-surfaces.

We define a local orthonormal basis $\left( \widehat{\psi },\widehat{s},%
\widehat{\phi }\right) $ as 
\begin{eqnarray}
\widehat{\psi } &=&\frac{\nabla \psi }{\left| \nabla \psi \right| } \\
\widehat{s} &=&\widehat{\phi }\times \widehat{\psi }
\end{eqnarray}

The transformation from $\left( \widehat{r},\widehat{\theta }\right) $ to $%
\left( \widehat{\psi },\widehat{s}\right) $ is a rotation of angle $\alpha $
such that 
\begin{equation}
\left( 
\begin{array}{l}
\widehat{\psi } \\ 
\widehat{s}
\end{array}
\right) =\left( 
\begin{array}{ll}
\cos \alpha & -\sin \alpha \\ 
\sin \alpha & \cos \alpha
\end{array}
\right) \cdot \left( 
\begin{array}{l}
\widehat{r} \\ 
\widehat{\theta }
\end{array}
\right)
\end{equation}

\subsubsection{Position Vector}

The position vector remains 
\begin{equation}
{\bf X}=R_{0}\widehat{R}+Z_{0}\widehat{Z}+r\widehat{r}
\end{equation}

\subsubsection{Covariant Basis}

The covariant vector basis is defined in (\ref{cov}), which becomes here 
\begin{eqnarray}
{\bf e}_{\psi } &=&\frac{\partial {\bf X}}{\partial \psi }=\left. \frac{%
\partial r}{\partial \psi }\right| _{s}\widehat{r}+r\left. \frac{\partial 
\widehat{r}}{\partial \psi }\right| _{s}=\left. \frac{\partial r}{\partial
\psi }\right| _{s}\widehat{r}+r\left. \frac{\partial \theta }{\partial \psi }%
\right| _{s}\widehat{\theta } \\
{\bf e}_{s} &=&\frac{\partial {\bf X}}{\partial s}=\left. \frac{\partial r}{%
\partial s}\right| _{\psi }\widehat{r}+r\left. \frac{\partial \widehat{r}}{%
\partial s}\right| _{\psi }=\left. \frac{\partial r}{\partial s}\right|
_{\psi }\widehat{r}+r\left. \frac{\partial \theta }{\partial s}\right|
_{\psi }\widehat{\theta } \\
{\bf e}_{\phi } &=&\frac{\partial {\bf X}}{\partial \phi }=R_{0}\frac{%
\partial \widehat{R}}{\partial \phi }+r\frac{\partial \widehat{r}}{\partial
\phi }=\left( R_{0}+r\cos \theta \right) \frac{\partial \widehat{R}}{%
\partial \phi }=R\widehat{\phi }
\end{eqnarray}
so that we have the covariant basis 
\begin{equation}
\left( {\bf e}_{\psi },{\bf e}_{s},{\bf e}_{\phi }\right) =\left( \left. 
\frac{\partial r}{\partial \psi }\right| _{s}\widehat{r}+r\left. \frac{%
\partial \theta }{\partial \psi }\right| _{s}\widehat{\theta },\left. \frac{%
\partial r}{\partial s}\right| _{\psi }\widehat{r}+r\left. \frac{\partial
\theta }{\partial s}\right| _{\psi }\widehat{\theta },R\widehat{\phi }\right)
\label{covpsis}
\end{equation}
the scaling factors 
\begin{equation}
\left( h_{\psi },h_{s},h_{\phi }\right) =\left( \sqrt{\left. \frac{\partial r%
}{\partial \psi }\right| _{s}^{2}+r^{2}\left. \frac{\partial \theta }{%
\partial \psi }\right| _{s}^{2}},\sqrt{\left. \frac{\partial r}{\partial s}%
\right| _{\psi }^{2}+r^{2}\left. \frac{\partial \theta }{\partial s}\right|
_{\psi }^{2}},R\right)
\end{equation}
and the normalized tangent basis 
\begin{equation}
\left( \widehat{e}_{\psi },\widehat{e}_{s},\widehat{e}_{\phi }\right)
=\left( \frac{1}{h_{\psi }}\left[ \left. \frac{\partial r}{\partial \psi }%
\right| _{s}\widehat{r}+r\left. \frac{\partial \theta }{\partial \psi }%
\right| _{s}\widehat{\theta }\right] ,\frac{1}{h_{s}}\left[ \left. \frac{%
\partial r}{\partial s}\right| _{\psi }\widehat{r}+r\left. \frac{\partial
\theta }{\partial s}\right| _{\psi }\widehat{\theta }\right] ,\widehat{\phi }%
\right)
\end{equation}

\subsubsection{Contravariant Basis}

The Contravariant vector basis is defined in (\ref{cont}), which becomes
here 
\begin{eqnarray}
{\bf e}^{\psi } &=&\nabla \psi =\left| \nabla \psi \right| \widehat{\psi } \\
{\bf e}^{s} &=&\nabla s=\widehat{s} \\
{\bf e}^{\phi } &=&\nabla \phi =\frac{\widehat{\phi }}{R}
\end{eqnarray}

The relations (\ref{rel2}) then give 
\begin{eqnarray}
{\bf e}_{\psi } &=&\frac{{\bf e}^{s}\times {\bf e}^{\phi }}{{\bf e}^{\psi
}\cdot {\bf e}^{s}\times {\bf e}^{\phi }}=\frac{\widehat{\psi }}{\left|
\nabla \psi \right| }  \label{covpsis2} \\
{\bf e}_{s} &=&\frac{{\bf e}^{\phi }\times {\bf e}^{\psi }}{{\bf e}^{s}\cdot 
{\bf e}^{\phi }\times {\bf e}^{\psi }}=\widehat{s}  \nonumber \\
{\bf e}_{\phi } &=&\frac{{\bf e}^{\psi }\times {\bf e}^{s}}{{\bf e}^{\phi
}\cdot {\bf e}^{\psi }\times {\bf e}^{s}}=R\widehat{\phi }  \nonumber
\end{eqnarray}
so that we have the following tangent basis 
\begin{equation}
\left( {\bf e}_{\psi },{\bf e}_{s},{\bf e}_{\phi }\right) =\left( \frac{%
\widehat{\psi }}{\left| \nabla \psi \right| },\widehat{s},R\widehat{\phi }%
\right)
\end{equation}
the scaling factors 
\begin{equation}
\left( h_{\psi },h_{s},h_{\phi }\right) =\left( \frac{1}{\left| \nabla \psi
\right| },1,R\right)
\end{equation}
the normalized tangent basis 
\begin{equation}
\left( \widehat{e}_{\psi },\widehat{e}_{s},\widehat{e}_{\phi }\right)
=\left( \widehat{\psi },\widehat{s},\widehat{\phi }\right)
\end{equation}
the reciprocal basis 
\begin{equation}
\left( {\bf e}^{\psi },{\bf e}^{s},{\bf e}^{\phi }\right) =\left( \left|
\nabla \psi \right| \widehat{\psi },\widehat{s},\frac{\widehat{\phi }}{R}%
\right)
\end{equation}
and the normalized reciprocal basis 
\begin{equation}
\left( \widehat{e}^{\psi },\widehat{e}^{s},\widehat{e}^{\phi }\right)
=\left( \widehat{\psi },\widehat{s},\widehat{\phi }\right)
\end{equation}
which here coincides with the normalized tangent basis, since both bases are
orthogonal.

By comparing (\ref{covpsis}) with (\ref{covpsis2}), we also find that 
\begin{eqnarray}
\left. \frac{\partial r}{\partial \psi }\right| _{s} &=&\frac{\cos \alpha }{%
\left| \nabla \psi \right| } \\
\left. \frac{\partial \theta }{\partial \psi }\right| _{s} &=&\frac{-\sin
\alpha }{r\left| \nabla \psi \right| } \\
\left. \frac{\partial r}{\partial s}\right| _{\psi } &=&\left( \widehat{s}%
\cdot \widehat{r}\right) =\sin \alpha \\
\left. \frac{\partial \theta }{\partial s}\right| _{\psi } &=&\frac{\left( 
\widehat{s}\cdot \widehat{\theta }\right) }{r}=\frac{\cos \alpha }{r}
\end{eqnarray}

\subsubsection{Metric Coefficients}

They are defined in (\ref{met}) and become here 
\begin{eqnarray}
g_{ij} &=&\left( 
\begin{array}{lll}
1/\left| \nabla \psi \right| ^{2} & 0 & 0 \\ 
0 & 1 & 0 \\ 
0 & 0 & R^{2}
\end{array}
\right) \\
g^{ij} &=&\left( 
\begin{array}{lll}
\left| \nabla \psi \right| ^{2} & 0 & 0 \\ 
0 & 1 & 0 \\ 
0 & 0 & 1/R^{2}
\end{array}
\right)  \nonumber
\end{eqnarray}

As a result 
\begin{equation}
g=\frac{R^{2}}{\left| \nabla \psi \right| ^{2}}
\end{equation}
and the Jacobian is 
\begin{equation}
J=\frac{R}{\left| \nabla \psi \right| }
\end{equation}

\subsubsection{Differential elements}

\begin{eqnarray}
dl\left( \psi \right) &=&\frac{d\psi }{\left| \nabla \psi \right| } \\
dl\left( s\right) &=&ds \\
dl\left( \phi \right) &=&Rd\phi
\end{eqnarray}

\begin{eqnarray}
d{\bf S}\left( \psi \right) &=&Rdsd\phi \widehat{\psi } \\
d{\bf S}\left( s\right) &=&\frac{R}{\left| \nabla \psi \right| }d\psi d\phi 
\widehat{s} \\
d{\bf S}\left( \phi \right) &=&\frac{1}{\left| \nabla \psi \right| }d\psi ds%
\widehat{\phi }
\end{eqnarray}

\begin{equation}
d^{3}{\bf X}=\frac{R}{\left| \nabla \psi \right| }d\psi dsd\phi
\end{equation}

\subsubsection{Christoffel Symbols}

They are defined in (\ref{chris}) and are here

\subsubsection{Differential Operations}

\paragraph{Gradient}

\begin{equation}
\nabla f=\left| \nabla \psi \right| \frac{\partial f}{\partial \psi }%
\widehat{\psi }+\frac{\partial f}{\partial s}\widehat{s}+\frac{1}{R}\frac{%
\partial f}{\partial \phi }\widehat{\phi }
\end{equation}

\paragraph{Divergence}

\begin{equation}
\nabla \cdot {\bf A}=\frac{\left| \nabla \psi \right| }{R}\frac{\partial }{%
\partial \psi }\left( R{\bf A}\cdot \widehat{\psi }\right) +\frac{\left|
\nabla \psi \right| }{R}\frac{\partial }{\partial s}\left( \frac{R}{\left|
\nabla \psi \right| }{\bf A}\cdot \widehat{s}\right) +\frac{1}{R}\frac{%
\partial }{\partial \phi }\left( {\bf A}\cdot \widehat{\phi }\right)
\end{equation}

\paragraph{Curl}

\begin{eqnarray}
\left( \nabla \times {\bf A}\right) \cdot \widehat{\psi } &=&\frac{1}{R}%
\frac{\partial }{\partial s}\left( R{\bf A}\cdot \widehat{\phi }\right) -%
\frac{1}{R}\frac{\partial }{\partial \phi }\left( {\bf A}\cdot \widehat{s}%
\right) \\
\left( \nabla \times {\bf A}\right) \cdot \widehat{s} &=&\frac{1}{R}\frac{%
\partial }{\partial \phi }\left( {\bf A}\cdot \widehat{\psi }\right) -\frac{%
\left| \nabla \psi \right| }{R}\frac{\partial }{\partial \psi }\left( R{\bf A%
}\cdot \widehat{\phi }\right) \\
\left( \nabla \times {\bf A}\right) \cdot \widehat{\phi } &=&\left| \nabla
\psi \right| \frac{\partial }{\partial \psi }\left( {\bf A}\cdot \widehat{s}%
\right) -\left| \nabla \psi \right| \frac{\partial }{\partial s}\left( \frac{%
{\bf A}\cdot \widehat{\psi }}{\left| \nabla \psi \right| }\right)
\end{eqnarray}

\subsection{System $\left( \psi ,\theta ,\phi \right) $}

\subsubsection{Definition}

The coordinates $\left( \psi ,\theta ,\phi \right) $ are defined from the
origin $\left( R_{0},Z_{0}\right) $ on the space 
\[
\min \left( \psi _{0},\psi _{a}\right) \leq \psi <\max \left( \psi _{0},\psi
_{a}\right) 
\]
and is related to $\left( r,\theta ,\phi \right) $ by

\[
\psi =\psi \left( r,\theta \right) 
\]
which is inverted to 
\[
r=r\left( \psi ,\theta \right) 
\]

Note that $\psi \left( r,\theta \right) $ must be a monotonic function of $r$
from $\psi _{0}$ at the center $\left( R_{0},Z_{0}\right) $ to $\psi _{a}$
at the edge. It is the case for nested flux-surfaces.

\subsubsection{Position Vector}

The position vector then becomes 
\begin{equation}
{\bf X}=R_{0}\widehat{R}+Z_{0}\widehat{Z}+r\left( \psi ,\theta \right) 
\widehat{r}
\end{equation}

\subsubsection{Covariant Basis}

The covariant vector basis is defined in (\ref{cov}), which becomes here 
\begin{eqnarray}
{\bf e}_{\psi } &=&\frac{\partial {\bf X}}{\partial \psi }=\left. \frac{%
\partial r}{\partial \psi }\right| _{\theta }\widehat{r} \\
{\bf e}_{\theta } &=&\frac{\partial {\bf X}}{\partial \theta }=\left. \frac{%
\partial r}{\partial \theta }\right| _{\psi }\widehat{r}+r\frac{\partial 
\widehat{r}}{\partial \theta }=\left. \frac{\partial r}{\partial \theta }%
\right| _{\psi }\widehat{r}+r\widehat{\theta } \\
{\bf e}_{\phi } &=&\frac{\partial {\bf X}}{\partial \phi }=R_{0}\frac{%
\partial \widehat{R}}{\partial \phi }+r\frac{\partial \widehat{r}}{\partial
\phi }=\left( R_{0}+r\cos \theta \right) \frac{\partial \widehat{R}}{%
\partial \phi }=R\widehat{\phi }
\end{eqnarray}
so that we have the covariant basis 
\begin{equation}
\left( {\bf e}_{\psi },{\bf e}_{\theta },{\bf e}_{\phi }\right) =\left(
\left. \frac{\partial r}{\partial \psi }\right| _{\theta }\widehat{r},\left. 
\frac{\partial r}{\partial \theta }\right| _{\psi }\widehat{r}+r\widehat{%
\theta },R\widehat{\phi }\right)  \label{covpsit}
\end{equation}
the scaling factors 
\begin{equation}
\left( h_{\psi },h_{\theta },h_{\phi }\right) =\left( \left| \left. \frac{%
\partial r}{\partial \psi }\right| _{\theta }\right| ,\sqrt{\left. \frac{%
\partial r}{\partial \theta }\right| _{\psi }^{2}+r^{2}},R\right)
\end{equation}
and the normalized tangent basis 
\begin{equation}
\left( \widehat{e}_{\psi },\widehat{e}_{\theta },\widehat{e}_{\phi }\right)
=\left( \widehat{r},\frac{1}{h_{\theta }}\left[ \left. \frac{\partial r}{%
\partial \theta }\right| _{\psi }\widehat{r}+r\widehat{\theta }\right] ,%
\widehat{\phi }\right)
\end{equation}

\subsubsection{Contravariant Basis}

The Contravariant vector basis is defined in (\ref{cont}), which becomes
here 
\begin{eqnarray}
{\bf e}^{\psi } &=&\nabla \psi =\left| \nabla \psi \right| \widehat{\psi } \\
{\bf e}^{\theta } &=&\nabla \theta =\frac{\widehat{\theta }}{r} \\
{\bf e}^{\phi } &=&\nabla \phi =\frac{\widehat{\phi }}{R}
\end{eqnarray}

The relations (\ref{rel2}) then give 
\begin{eqnarray}
{\bf e}_{\psi } &=&\frac{{\bf e}^{\theta }\times {\bf e}^{\phi }}{{\bf e}%
^{\psi }\cdot {\bf e}^{\theta }\times {\bf e}^{\phi }}=\frac{\widehat{r}}{%
\left| \nabla \psi \right| \cos \alpha } \\
{\bf e}_{\theta } &=&\frac{{\bf e}^{\phi }\times {\bf e}^{\psi }}{{\bf e}%
^{\theta }\cdot {\bf e}^{\phi }\times {\bf e}^{\psi }}=\frac{r\widehat{s}}{%
\left( \widehat{\theta }\cdot \widehat{s}\right) }=\frac{r\widehat{s}}{\cos
\alpha } \\
{\bf e}_{\phi } &=&\frac{{\bf e}^{\psi }\times {\bf e}^{\theta }}{{\bf e}%
^{\phi }\cdot {\bf e}^{\psi }\times {\bf e}^{\theta }}=R\widehat{\phi }
\end{eqnarray}
so that we have the following tangent basis 
\begin{equation}
\left( {\bf e}_{\psi },{\bf e}_{\theta },{\bf e}_{\phi }\right) =\left( 
\frac{\widehat{r}}{\left| \nabla \psi \right| \cos \alpha },\frac{r\widehat{s%
}}{\cos \alpha },R\widehat{\phi }\right)  \label{covpsit2}
\end{equation}
the scaling factors 
\begin{equation}
\left( h_{\psi },h_{\theta },h_{\phi }\right) =\left( \frac{1}{\left| \nabla
\psi \right| \cos \alpha },\frac{r}{\cos \alpha },R\right)
\end{equation}
the normalized tangent basis 
\begin{equation}
\left( \widehat{e}_{\psi },\widehat{e}_{\theta },\widehat{e}_{\phi }\right)
=\left( \widehat{r},\widehat{s},\widehat{\phi }\right)
\end{equation}
the reciprocal basis 
\begin{equation}
\left( {\bf e}^{\psi },{\bf e}^{\theta },{\bf e}^{\phi }\right) =\left(
\left| \nabla \psi \right| \widehat{\psi },\frac{\widehat{\theta }}{r},\frac{%
\widehat{\phi }}{R}\right)
\end{equation}
and the normalized reciprocal basis 
\begin{equation}
\left( \widehat{e}^{\psi },\widehat{e}^{\theta },\widehat{e}^{\phi }\right)
=\left( \widehat{\psi },\widehat{\theta },\widehat{\phi }\right)
\end{equation}
which here does not coincide with the normalized tangent basis, since both
bases are not orthogonal.

By comparing (\ref{covpsit}) with (\ref{covpsit2}), we also find that 
\begin{eqnarray}
\left. \frac{\partial r}{\partial \psi }\right| _{\theta } &=&\frac{1}{%
\left| \nabla \psi \right| \cos \alpha } \\
\left. \frac{\partial r}{\partial \theta }\right| _{\psi } &=&r\sqrt{\frac{1%
}{\left( \widehat{\theta }\cdot \widehat{s}\right) ^{2}}-1}=r\tan \alpha
\end{eqnarray}

\subsubsection{Metric Coefficients}

They are defined in (\ref{met}) and become here 
\begin{equation}
g_{ij}=\left( 
\begin{array}{lll}
1/\left[ \left| \nabla \psi \right| \cos \alpha \right] ^{2} & r\tan \alpha
/\left[ \left| \nabla \psi \right| \cos \alpha \right]  & 0 \\ 
r\tan \alpha /\left[ \left| \nabla \psi \right| \cos \alpha \right]  & 
r^{2}/\cos ^{2}\alpha  & 0 \\ 
0 & 0 & R^{2}
\end{array}
\right)   \label{gijpsitheta}
\end{equation}
or equivalently 
\begin{equation}
g_{ij}=\left( 
\begin{array}{lll}
\left. \partial r/\partial \psi \right| _{\theta }^{2} & \left. \partial r/%
\partial \psi \right| _{\theta }\left. \partial r/\partial \theta \right|
_{\psi } & 0 \\ 
\left. \partial r/\partial \psi \right| _{\theta }\left. \partial r/\partial %
\theta \right| _{\psi } & \left. \partial r/\partial \theta \right| _{\psi
}^{2}+r^{2} & 0 \\ 
0 & 0 & R^{2}
\end{array}
\right)   \nonumber
\end{equation}
and 
\begin{equation}
g^{ij}=\left( 
\begin{array}{lll}
\left| \nabla \psi \right| ^{2} & -\left| \nabla \psi \right| \sin \alpha /r
& 0 \\ 
-\left| \nabla \psi \right| \sin \alpha /r & 1/r^{2} & 0 \\ 
0 & 0 & 1/R^{2}
\end{array}
\right) 
\end{equation}
As a result 
\begin{equation}
g=\frac{R^{2}r^{2}}{\left| \nabla \psi \right| ^{2}\cos ^{2}\alpha }
\end{equation}
and the Jacobian is 
\begin{equation}
J=\frac{Rr}{\left| \nabla \psi \right| \cos \alpha }  \label{jacpsitheta}
\end{equation}

\subsubsection{Differential elements}

\begin{eqnarray}
dl\left( \psi \right) &=&\frac{d\psi }{\left| \nabla \psi \right| \cos
\alpha } \\
dl\left( \theta \right) &=&\frac{r}{\cos \alpha }d\theta \\
dl\left( \phi \right) &=&Rd\phi
\end{eqnarray}

\begin{eqnarray}
d{\bf S}\left( \psi \right) &=&\frac{Rr}{\cos \alpha }d\theta d\phi \widehat{%
\psi } \\
d{\bf S}\left( \theta \right) &=&\frac{R}{\left| \nabla \psi \right| \cos
\alpha }d\psi d\phi \widehat{s} \\
d{\bf S}\left( \phi \right) &=&\frac{r}{\left| \nabla \psi \right| \cos
\alpha }d\psi d\theta \widehat{\phi }
\end{eqnarray}

\begin{equation}
d^{3}{\bf X}=\frac{Rr}{\left| \nabla \psi \right| \cos \alpha }d\psi d\theta
d\phi
\end{equation}

\subsubsection{Christoffel Symbols}

They are defined in (\ref{chris}) and are here

\subsubsection{Differential Operations}

\paragraph{Gradient}

\begin{equation}
\nabla f=\left| \nabla \psi \right| \frac{\partial f}{\partial \psi }%
\widehat{\psi }+\frac{1}{r}\frac{\partial f}{\partial \theta }\widehat{%
\theta }+\frac{1}{R}\frac{\partial f}{\partial \phi }\widehat{\phi }
\end{equation}

\paragraph{Divergence}

\begin{equation}
\nabla \cdot {\bf A}=\frac{\left| \nabla \psi \right| \cos \alpha }{Rr}\frac{%
\partial }{\partial \psi }\left( \frac{Rr}{\cos \alpha }{\bf A}\cdot 
\widehat{\psi }\right) +\frac{\left| \nabla \psi \right| \cos \alpha }{Rr}%
\frac{\partial }{\partial \theta }\left( \frac{R}{\left| \nabla \psi \right|
\cos \alpha }{\bf A}\cdot \widehat{\theta }\right) +\frac{1}{R}\frac{%
\partial }{\partial \phi }\left( {\bf A}\cdot \widehat{\phi }\right)
\end{equation}

\paragraph{Curl}

\begin{eqnarray}
\left( \nabla \times {\bf A}\right) \cdot \widehat{\psi } &=&\frac{\cos
\alpha }{Rr}\frac{\partial }{\partial \theta }\left( R{\bf A}\cdot \widehat{%
\phi }\right) -\frac{1}{R}\frac{\partial }{\partial \phi }\left( {\bf A}%
\cdot \widehat{s}\right) \\
\left( \nabla \times {\bf A}\right) \cdot \widehat{\theta } &=&\frac{1}{R}%
\frac{\partial }{\partial \phi }\left( {\bf A}\cdot \widehat{r}\right) -%
\frac{\left| \nabla \psi \right| \cos \alpha }{R}\frac{\partial }{\partial
\psi }\left( R{\bf A}\cdot \widehat{\phi }\right) \\
\left( \nabla \times {\bf A}\right) \cdot \widehat{\phi } &=&\frac{\left|
\nabla \psi \right| \cos \alpha }{r}\frac{\partial }{\partial \psi }\left( 
\frac{r{\bf A}\cdot \widehat{s}}{\cos \alpha }\right) -\frac{\left| \nabla
\psi \right| \cos \alpha }{r}\frac{\partial }{\partial \theta }\left( \frac{%
{\bf A}\cdot \widehat{r}}{\left| \nabla \psi \right| \cos \alpha }\right)
\end{eqnarray}
